\section{Hoja de ruta con casos: sprint de 12 semanas de investigación abierta en Agent}

\subsection{De la filosofía del curso al plan de ejecución}
Para convertir la filosofía de esta clase en un proyecto ejecutable, se puede adoptar una estructura de «sprint de 12 semanas»: las primeras 4 semanas se dedican a construir la infraestructura, las 4 semanas intermedias a la iteración del sistema, y las últimas 4 semanas a la robustez y al empaquetado para la reproducción. Este ritmo corresponde al espíritu central del expositor: \textbf{primero convertir el sistema de investigación en una máquina verificable, y solo después buscar picos de desempeño locales.}

\begin{figure}[H]
\centering
\includegraphics[width=0.93\textwidth]{frame-02-benefits.jpg}
\caption{Objetivo de beneficios del sistema Agent: la trinidad de interfaz, capacidad y arquitectura}
\end{figure}
\noindent\footnotesize Evidencia visual: los subtítulos se ubican aproximadamente entre 00:00:43 y 00:01:39; el expositor extiende «access shapes research» hacia los límites de capacidad del sistema.\normalsize

\begin{importantbox}{El objetivo sugerido del sprint no es «lanzar una demo llamativa»}
Un objetivo de mayor valor es:
\begin{itemize}
    \item que cualquier conclusión experimental pueda ser reproducida por un compañero en 24 horas;
    \item que cualquier trayectoria de fallo pueda ubicarse rápidamente en una etapa concreta;
    \item que cualquier mejora pueda explicarse como proveniente de los datos, la estrategia o la cadena de herramientas.
\end{itemize}
Hacer esto parece lento en el corto plazo, pero a largo plazo reduce de manera significativa la proporción de «progreso falso».
\end{importantbox}

\subsection{Hitos por etapa y entregables}
\begin{table}[H]
\centering
\begin{tabular}{p{2.2cm}p{3.5cm}p{3.6cm}p{3.2cm}}
\toprule
\textbf{Etapa} & \textbf{Tarea central} & \textbf{Entregable obligatorio} & \textbf{Criterio de aceptación} \\
\midrule
Semanas 1--2 & Unificar el protocolo de tareas y el formato de registro & task schema, run schema, replay script & Cualquier tarea puede reproducirse \\
Semanas 3--4 & Hacer funcionar el agent de línea base y la evaluación & baseline result, failure taxonomy & Existe una varianza de línea base estable \\
Semanas 5--6 & Incorporar los módulos de memoria, reflexión y planificación & ablation report, trace dashboard & Se puede explicar el origen de la ganancia \\
Semanas 7--8 & Orquestación de múltiples agent y gobernanza de permisos & orchestration spec, safety gate & Las conductas que exceden los permisos pueden bloquearse y auditarse \\
Semanas 9--10 & Robustez y pruebas de estrés fuera de distribución & stress suite, adversarial suite & Los fallos bajo distintas distribuciones pueden descomponerse \\
Semanas 11--12 & Empaquetado para la reproducción y publicación documentada & reproducibility package, model card & Terceros pueden reproducirlo de manera independiente \\
\bottomrule
\end{tabular}
\caption{Ruta mínima de ingeniería de investigación del sprint de 12 semanas}
\end{table}

\begin{knowledgebox}{Por qué hacer primero el replay y después la optimización}
Muchos equipos dedican el 80\% del tiempo a «pensar nuevas estrategias» sin contar con trayectorias reproducibles.  
Sin replay, no se puede determinar si una mejora se debe a que la estrategia funciona o a una fluctuación aleatoria; sin un replay estable, ninguna optimización posterior puede acumularse como conocimiento confiable.
\end{knowledgebox}

\subsection{Ciclo principal experimental sugerido (pseudocódigo)}
\begin{lstlisting}[caption={Research Loop: ciclo principal de iteración reproducible de Agent},label={lst:research-loop}]
for task_batch in benchmark:
    run_id = launch_agent(task_batch, config, tools, safety_policy)
    traces = collect_traces(run_id)
    metrics = evaluate(traces, eval_protocol)

    if metrics.regression_detected:
        rollback(config)
        tag_failure(traces, taxonomy)
        continue

    insights = reflection_and_ablation(traces, metrics)
    config = update_config(config, insights)
    snapshot(run_id, config, metrics, traces)
\end{lstlisting}

\begin{warningbox}{La deuda técnica más fácil de pasar por alto: experimentos no comparables}
Si en cada iteración se cambian el conjunto de tareas, la versión de las herramientas o el script de evaluación, surge la ilusión de que «siempre se está progresando, pero no hay dos resultados que puedan compararse».  
Durante la etapa del sprint es necesario congelar los protocolos clave y ajustar las variables solo dentro de una ventana controlada.
\end{warningbox}

\subsection{Resumen de la sección}
El objetivo de esta hoja de ruta de 12 semanas es convertir la filosofía de «ciencia abierta + ingeniería de Agent» en un sistema ejecutable: reproducible, auditable e interpretable. Si falta cualquiera de los tres, el equipo tendrá que rehacer el trabajo repetidamente en la etapa de escalamiento.


